\documentclass[UTF8,11pt]{ctexart}
\usepackage[a4paper,margin=24mm]{geometry}
\usepackage{amsmath,amssymb,amsthm,mathtools,mathrsfs}
\usepackage[all]{xy}
\usepackage{xurl,hyperref,enumitem}
\usepackage{tikz}
\usetikzlibrary{arrows.meta,cd,decorations.markings,calc,patterns}
\hypersetup{hidelinks,breaklinks=true}
\setlength{\emergencystretch}{3em}
\allowdisplaybreaks
\newtheorem*{theorem}{Theorem}
\newtheorem*{thm}{Theorem}
\newtheorem*{lemma}{Lemma}
\newtheorem*{proposition}{Proposition}
\newtheorem*{prop}{Proposition}
\newtheorem*{corollary}{Corollary}
\newtheorem*{cor}{Corollary}
\newtheorem*{claim}{Claim}
\theoremstyle{definition}
\newtheorem*{definition}{Definition}
\newtheorem*{defn}{Definition}
\newtheorem*{remark}{Remark}
\newtheorem*{example}{Example}
\newcommand{\R}{\mathbb R}
\newcommand{\C}{\mathbb C}
\newcommand{\Z}{\mathbb Z}
\newcommand{\Q}{\mathbb Q}
\newcommand{\N}{\mathbb N}
\newcommand{\F}{\mathbb F}
\newcommand{\D}{\mathbb D}
\newcommand{\bB}{\mathbb B}
\newcommand{\bC}{\mathbb C}
\newcommand{\bR}{\mathbb R}
\newcommand{\bZ}{\mathbb Z}
\newcommand{\bN}{\mathbb N}
\newcommand{\bQ}{\mathbb Q}
\newcommand{\bD}{\mathbb D}
\newcommand{\bF}{\mathbb F}
\newcommand{\CP}{\mathbb{CP}}
\newcommand{\RP}{\mathbb{RP}}
\newcommand{\sO}{\mathscr O}
\newcommand{\sI}{\mathscr I}
\newcommand{\sF}{\mathscr F}
\newcommand{\sG}{\mathscr G}
\newcommand{\sH}{\mathscr H}
\newcommand{\sR}{\mathscr R}
\newcommand{\sM}{\mathscr M}
\newcommand{\sV}{\mathscr V}
\newcommand{\sU}{\mathscr U}
\DeclareMathOperator{\Hom}{Hom}
\DeclareMathOperator{\Ext}{Ext}
\DeclareMathOperator{\Tor}{Tor}
\DeclareMathOperator{\Spec}{Spec}
\DeclareMathOperator{\Proj}{Proj}
\DeclareMathOperator{\supp}{supp}
\DeclareMathOperator{\im}{im}
\DeclareMathOperator{\id}{id}
\DeclareMathOperator{\Tr}{Tr}
\DeclareMathOperator{\tr}{tr}
\DeclareMathOperator{\rank}{rank}
\DeclareMathOperator{\ord}{ord}
\DeclareMathOperator{\dist}{dist}
\DeclareMathOperator{\End}{End}
\DeclareMathOperator{\Aut}{Aut}
\DeclareMathOperator{\Gal}{Gal}
\DeclareMathOperator{\vol}{vol}
\DeclareMathOperator{\Cl}{Cl}
\DeclareMathOperator{\coker}{coker}
\DeclareMathOperator{\codim}{codim}
\DeclareMathOperator{\divergence}{div}
\DeclareMathOperator{\Ric}{Ric}
\DeclareMathOperator{\grad}{grad}
\DeclareMathOperator{\Hess}{Hess}
\newcommand{\abs}[1]{\left\lvert#1\right\rvert}
\newcommand{\sabs}[1]{\lvert#1\rvert}
\newcommand{\babs}[1]{\bigl\lvert#1\bigr\rvert}
\newcommand{\Babs}[1]{\Bigl\lvert#1\Bigr\rvert}
\newcommand{\bbabs}[1]{\biggl\lvert#1\biggr\rvert}
\newcommand{\BBabs}[1]{\Biggl\lvert#1\Biggr\rvert}
\newcommand{\norm}[1]{\left\lVert#1\right\rVert}
\newcommand{\snorm}[1]{\lVert#1\rVert}
\newcommand{\bnorm}[1]{\bigl\lVert#1\bigr\rVert}
\newcommand{\Bnorm}[1]{\Bigl\lVert#1\Bigr\rVert}
\newcommand{\bbnorm}[1]{\biggl\lVert#1\biggr\rVert}
\newcommand{\BBnorm}[1]{\Biggl\lVert#1\Biggr\rVert}
\newcommand{\widebar}[1]{\overline{#1}}
\newcommand{\nicefrac}[2]{\frac{#1}{#2}}
\newcommand{\linnprod}[2]{\langle#1,#2\rangle}
\newcommand{\myindex}[1]{#1}
\newcommand{\glsadd}[1]{}
\newcommand{\avoidbreak}{}
\newcommand{\sourceref}[1]{\textnormal{[原文：\nolinkurl{#1}]}}
\newcommand{\thmref}[1]{\sourceref{#1}}
\newcommand{\lemmaref}[1]{\sourceref{#1}}
\newcommand{\propref}[1]{\sourceref{#1}}
\newcommand{\corref}[1]{\sourceref{#1}}
\newcommand{\defnref}[1]{\sourceref{#1}}
\newcommand{\exampleref}[1]{\sourceref{#1}}
\newcommand{\exerciseref}[1]{\sourceref{#1}}
\newcommand{\figureref}[1]{\sourceref{#1}}
\newcommand{\sectionref}[1]{\sourceref{#1}}
\newcommand{\appendixref}[1]{\sourceref{#1}}
\newcommand{\stacksref}[2]{\href{https://stacks.math.columbia.edu/tag/#1}{#2 [#1]}}


\newcommand{\E}{\mathbb E}
\newcommand{\Prob}{\mathbb P}
\newcommand{\Reff}{\mathcal R_{\mathrm{eff}}}
\newcommand{\eps}{\varepsilon}
\DeclareMathOperator{\diam}{diam}
\DeclareMathOperator{\Var}{Var}
\DeclareMathOperator{\Crit}{Crit}
\newcommand{\dd}{\,\mathrm d}

\begin{document}
\section*{067\quad 有限呈示代数上模的泛自由性}
\subsection*{来源与计数目标}
The Stacks Project Authors, \emph{The Stacks Project}, Lemma 10.118.2；支撑 Lemmas 10.118.1、10.115.7 (07NA)。在线正文访问日期2026-09-28。
\par\url{https://stacks.math.columbia.edu/tag/051Q}
\par\textbf{本文件只计一个问题：}整环 $R$ 上有限呈示代数 $S$ 的有限呈示模 $M$，在倒置一个非零底环元素后成为自由 $R_f$-模。
\subsection*{作者命题与人类证明}
\begin{lemma}[10.115.7: normalization after localization]
Let $R\to S$ be an injective finite type ring map. Assume $R$ is a domain. Then there exists an integer $d$ and a factorization
\[R\to R[y_1,\ldots,y_d]\to S'\to S\]
by injective maps such that $S'$ is finite over $R[y_1,\ldots,y_d]$ and such that $S'_f\cong S_f$ for some nonzero $f\in R$.
\end{lemma}
\begin{proof}
Pick $x_1,\ldots,x_n\in S$ which generate $S$ over $R$. Let $K$ be the fraction field of $R$ and $S_K=S\otimes_RK$. By Lemma 10.115.4 we can find $y_1,\ldots,y_d\in S$ such that $K[y_1,\ldots,y_d]\to S_K$ is a finite injective map. Note that $y_i\in S$ because we may pick the $y_j$ in the $\mathbf Z$-algebra generated by $x_1,\ldots,x_n$.

As a finite ring map is integral (see Lemma 10.36.3) we can find monic $P_i\in K[y_1,\ldots,y_d][T]$ such that $P_i(x_i)=0$ in $S_K$. Let $f\in R$ be a nonzero element such that $fP_i\in R[y_1,\ldots,y_d][T]$ for all $i$. Then $fP_i(x_i)$ maps to zero in $S_K$. Hence after replacing $f$ by another nonzero element of $R$ we may also assume $fP_i(x_i)$ is zero in $S$.

Set $x'_i=fx_i$ and let $S'\subset S$ be the $R$-subalgebra generated by $y_1,\ldots,y_d$ and $x'_1,\ldots,x'_n$. Note that $x'_i$ is integral over $R[y_1,\ldots,y_d]$ as we have $Q_i(x'_i)=0$ where $Q_i=f^{\deg_T(P_i)}P_i(T/f)$ which is a monic polynomial in $T$ with coefficients in $R[y_1,\ldots,y_d]$ by our choice of $f$. Hence $R[y_1,\ldots,y_d]\subset S'$ is finite by Lemma 10.36.5. Since $S'\subset S$ we have $S'_f\subset S_f$ (localization is exact).

On the other hand, the elements $x_i=x'_i/f$ in $S'_f$ generate $S_f$ over $R_f$ and hence $S'_f\to S_f$ is surjective. Whence $S'_f\cong S_f$ and we win.
\end{proof}

\begin{lemma}[10.118.1: generic freeness over a Noetherian domain]
Let $R\to S$ be a ring map. Let $M$ be an $S$-module. Assume $R$ is Noetherian, $R$ is a domain, $R\to S$ is of finite type, and $M$ is a finite type $S$-module. Then there exists a nonzero $f\in R$ such that $M_f$ is a free $R_f$-module.
\end{lemma}
\begin{proof}
Let $K$ be the fraction field of $R$. Set $S_K=K\otimes_RS$. This is an algebra of finite type over $K$. We will argue by induction on $d=\dim(S_K)$ (which is finite for example by Noether normalization, see Section 10.115). Fix $d\geq0$. Assume we know that the lemma holds in all cases where $\dim(S_K)<d$.

Suppose given $R\to S$ and $M$ as in the lemma with $\dim(S_K)=d$. By Lemma 10.62.1 there exists a filtration $0\subset M_1\subset M_2\subset\ldots\subset M_n=M$ so that $M_i/M_{i-1}$ is isomorphic to $S/\mathfrak q$ for some prime $\mathfrak q$ of $S$. Note that $\dim((S/\mathfrak q)_K)\leq\dim(S_K)$. Also, note that an extension of free modules is free (see basic notion 50). Thus we may assume $M=S$ and that $S$ is a domain of finite type over $R$.

If $R\to S$ has a nontrivial kernel, then take a nonzero $f\in R$ in this kernel. In this case $S_f=0$ and the lemma holds. (This is really the case $d=-1$ and the start of the induction.) Hence we may assume that $R\to S$ is a finite type extension of Noetherian domains.

Apply Lemma 10.115.7 and replace $R$ by $R_f$ (with $f$ as in the lemma) to get a factorization
\[R\subset R[y_1,\ldots,y_d]\subset S\]
where the second extension is finite. Choose $z_1,\ldots,z_r\in S$ which form a basis for the fraction field of $S$ over the fraction field of $R[y_1,\ldots,y_d]$. This gives a short exact sequence
\[0\to R[y_1,\ldots,y_d]^{\oplus r}\xrightarrow{(z_1,\ldots,z_r)}S\to N\to0.\]
By construction $N$ is a finite $R[y_1,\ldots,y_d]$-module whose support does not contain the generic point $(0)$ of $\Spec(R[y_1,\ldots,y_d])$. By Lemma 10.40.5 there exists a nonzero $g\in R[y_1,\ldots,y_d]$ such that $g$ annihilates $N$, so we may view $N$ as a finite module over $S'=R[y_1,\ldots,y_d]/(g)$. Since $\dim(S'_K)<d$ by induction there exists a nonzero $f\in R$ such that $N_f$ is a free $R_f$-module.

Since $(R[y_1,\ldots,y_d])_f\cong R_f[y_1,\ldots,y_d]$ is free also we conclude by the already mentioned fact that an extension of free modules is free.
\end{proof}

\begin{lemma}[10.118.2: generic freeness]
Let $R\to S$ be a ring map. Let $M$ be an $S$-module. Assume $R$ is a domain, $R\to S$ is of finite presentation, and $M$ is an $S$-module of finite presentation. Then there exists a nonzero $f\in R$ such that $M_f$ is a free $R_f$-module.
\end{lemma}
\begin{proof}
Write $S=R[x_1,\ldots,x_n]/(g_1,\ldots,g_m)$. For $g\in R[x_1,\ldots,x_n]$ denote $\overline g$ its image in $S$. We may write $M=S^{\oplus t}/\sum Sn_i$ for some $n_i\in S^{\oplus t}$. Write $n_i=(\overline g_{i1},\ldots,\overline g_{it})$ for some $g_{ij}\in R[x_1,\ldots,x_n]$. Let $R_0\subset R$ be the subring generated by all the coefficients of all the elements $g_i,g_{ij}\in R[x_1,\ldots,x_n]$.

Define $S_0=R_0[x_1,\ldots,x_n]/(g_1,\ldots,g_m)$. Define $M_0=S_0^{\oplus t}/\sum S_0n_i$. Then $R_0$ is a domain of finite type over $\mathbf Z$ and hence Noetherian (see Lemma 10.31.1). Moreover via the injection $R_0\to R$ we have $S\cong R\otimes_{R_0}S_0$ and $M\cong R\otimes_{R_0}M_0$. Applying Lemma 10.118.1 we obtain a nonzero $f\in R_0$ such that $(M_0)_f$ is a free $(R_0)_f$-module. Hence $M_f=R_f\otimes_{(R_0)_f}(M_0)_f$ is a free $R_f$-module.
\end{proof}
\subsection*{整理说明（非作者正文）}
本题主目标为10.118.2，不把10.118.1作为第二题，也不另计局部化正规化引理。标准先修为域上的 Noether 正规化（10.115.4，旧018材料可供阅读）、有限模的素商滤过、有限模支撑与零化子的关系、张量积/局部化基本性质，以及 Hilbert 基定理。

难度依据：先用一般纤维维数归纳，把模的自由子块与低维支撑余模分开，再通过有限系数子环将一般整环情形降至 Noetherian 情形。关键消分母、商模降维及有限系数下降均已收录。自由模可能具有无限秩，本条不附加有限秩结论。引用采用原编号，保存所用网页正文转录，未称作整章原件。
\subsection*{可修改的简短标注（非作者正文）}
$c$：有限呈示族中模的泛自由性。

$a$：Noether正规化；分式域基；素商滤过；支撑降维；自由扩张分裂；有限系数子环；基变换。

\texttt{cross\_theory\_status = yes}。将一般纤维的维数与模支撑转成可归纳的低维余模，再以系数子环和基变换把自由性带回原整环。位置：10.118.1的短正合列及10.118.2的有限系数构造。
标签为非穷尽、可修改初稿。
\subsection*{许可与修改记录}
原作者及作品见来源栏；来源采用 GNU Free Documentation License 1.2 或后续版本。本题证摘录和附加整理沿用该许可。2026-09-28：增加题号、来源定位、独立排版、整理说明和初稿标注；数学内容的取材范围及排版变化见整理说明。
\end{document}
